\providecommand{\SO}{\mathrm{SO}}
\providecommand{\SE}{\mathrm{SE}}
\providecommand{\Eucl}{\mathrm{E}}
\providecommand{\Hyp}{\mathbb{H}}
\providecommand{\Attn}{\operatorname{Attention}}
\providecommand{\softmax}{\operatorname{softmax}}
\providecommand{\Concat}{\operatorname{Concat}}
\providecommand{\sigm}{\sigma}
\providecommand{\Lc}{\mathcal{L}}

\chapter{Aplicaciones}
\label{ch:applications}

Los cap\'itulos anteriores construyeron una teor\'ia: los dominios poseen
simetr\'ias, las simetr\'ias restringen las arquitecturas y un \'unico
principio ---el paso de mensajes equivariante--- unifica las redes
convolucionales, las redes de grafos y los mecanismos de atenci\'on. Este
cap\'itulo se vuelca hacia la pr\'actica. Su prop\'osito no es catalogar todos
los sistemas que emplean aprendizaje profundo geom\'etrico, sino mostrar c\'omo
la maquinaria abstracta de las ``5~Gs'' se concreta en problemas reales y por
qu\'e el punto de vista geom\'etrico se traduce en ganancias medibles en
precisi\'on, eficiencia de datos y consistencia f\'isica.

El cap\'itulo tiene dos movimientos. El primero desarrolla el \emph{mecanismo de
atenci\'on} en profundidad: su historia, sus fundamentos algebraicos, la
geometr\'ia de sus espacios de consultas, claves y valores, su formulaci\'on
como operador integral y ---de manera crucial--- el sentido preciso en el que
constituye la contrapartida adaptativa de la convoluci\'on geom\'etrica de los
\Cref{ch:grids,ch:groups}. Este material es fundacional m\'as que meramente
aplicado: el \emph{transformer} es la arquitectura detr\'as de buena parte del
aprendizaje autom\'atico moderno, y comprenderlo geom\'etricamente explica tanto
su \'exito como la larga estirpe de variantes equivariantes que sustentan el
resto del cap\'itulo. El segundo movimiento recorre las propias \'areas de
aplicaci\'on, organizadas en torno a tres familias de datos. La \emph{visi\'on
por computador} abarca desde rejillas regulares de p\'ixeles hasta nubes de
puntos tridimensionales sin estructura, ejercitando los temas de las rejillas,
los grupos y los \emph{gauges}. Las \emph{mol\'eculas y los sistemas f\'isicos}
son el escaparate de la equivarianza $\Eucl(3)$ y $\SE(3)$, donde respetar las
simetr\'ias de la f\'isica no es un refinamiento sino un requisito. Los
\emph{sistemas de recomendaci\'on y las redes sociales} son relacionales por
naturaleza y se apoyan en la invarianza a permutaciones y, cada vez m\'as, en
geometr\'ias no euclidianas.

\section{Mecanismos de atenci\'on}
\label{sec:apps:attention-foundations}

Los mecanismos de atenci\'on son el complemento natural de las convoluciones
geom\'etricas desarrolladas en los cap\'itulos anteriores: donde las
convoluciones imponen una estructura r\'igida mediante \emph{kernels}
compartidos e independientes del contenido, la atenci\'on introduce flexibilidad
adaptativa, ponderando din\'amicamente la contribuci\'on de cada elemento seg\'un
la informaci\'on que transporta. Esta secci\'on desarrolla los fundamentos
matem\'aticos de la atenci\'on, desde su formulaci\'on algebraica cl\'asica hasta
las extensiones geom\'etricas que le permiten operar sobre dominios no
euclidianos preservando simetr\'ias.

\subsection{Historia y motivaci\'on}
\label{subsec:apps:attn-history}

Los mecanismos de atenci\'on surgieron hist\'oricamente de la necesidad de
superar las limitaciones de los modelos secuenciales en el procesamiento del
lenguaje natural. Los primeros trabajos de \citet{bahdanau2014neural} y
\citet{luong2015effective} introdujeron la atenci\'on en sistemas de traducci\'on
autom\'atica basados en redes recurrentes, permitiendo al modelo ``atender'' a
distintas partes de la secuencia de entrada durante la generaci\'on. El punto de
inflexi\'on lleg\'o con el trabajo seminal \emph{``Attention is All You Need''}
de \citet{vaswani2017attentionneed}, que demostr\'o que la atenci\'on por s\'i
sola ---sin recurrencia ni convoluci\'on--- pod\'ia alcanzar un rendimiento de
estado del arte en traducci\'on. Esta arquitectura \emph{transformer}
revolucion\'o no solo el procesamiento del lenguaje natural, sino que se
convirti\'o en el fundamento de modelos de visi\'on por computador, de
procesamiento de grafos y, finalmente, del aprendizaje profundo geom\'etrico.

La evoluci\'on hist\'orica puede resumirse en tres fases: la atenci\'on como
\emph{complemento} de las redes recurrentes (2014--2016), empleada para mejorar
el manejo de las dependencias a largo alcance; la atenci\'on como
\emph{arquitectura principal} (2017--2019), con \emph{transformers} puros que
reemplazan a los modelos recurrentes; y la \emph{atenci\'on geom\'etrica}
(2020--presente), que extiende el mecanismo a datos no euclidianos preservando
simetr\'ias. El resto del cap\'itulo sigue ese arco hasta su destino actual.

\subsection{Fundamentos algebraicos}
\label{subsec:apps:attn-algebra}

La atenci\'on se asienta sobre conceptos de \'algebra lineal y geometr\'ia de
espacios vectoriales. En esencia opera sobre tres espacios vectoriales
interrelacionados que codifican distintos aspectos de la informaci\'on.

\paragraph{Espacios de caracter\'isticas.}
Consideremos un espacio de caracter\'isticas $\mathcal{H} = \R^d$ en el que
residen las representaciones de entrada. Para un conjunto de $n$ elementos de
entrada los reunimos en una matriz
\[
  X = [x_1, x_2, \dots, x_n]^\top \in \R^{n \times d},
\]
donde cada $x_i \in \mathcal{H}$ es un elemento de la secuencia de entrada. El
mecanismo de atenci\'on define tres transformaciones lineales fundamentales, las
\emph{consultas}, las \emph{claves} y los \emph{valores},
\[
  Q = X W_Q \in \R^{n\times d_k},
  \qquad
  K = X W_K \in \R^{n\times d_k},
  \qquad
  V = X W_V \in \R^{n\times d_v},
\]
donde las matrices de peso $W_Q, W_K \in \R^{d\times d_k}$ y
$W_V \in \R^{d\times d_v}$ son par\'ametros entrenables que proyectan el espacio
de entrada sobre subespacios especializados. En la implementaci\'on multicabeza
est\'andar se toma $d_k = d_v = d/h$, donde $h$ es el n\'umero de cabezas de
atenci\'on.

\begin{definition}[El producto punto como medida de similitud]
\label{def:apps:dot-similarity}
La funci\'on de compatibilidad fundamental en atenci\'on es el producto punto
\[
  \operatorname{sim}(q_i, k_j) = q_i^\top k_j
  = \norm{q_i}\,\norm{k_j}\cos\theta_{ij},
\]
donde $\theta_{ij}$ es el \'angulo entre $q_i$ y $k_j$.
\end{definition}

Esta lectura geom\'etrica muestra que la atenci\'on mide una similitud coseno
normalizada entre representaciones en el espacio de caracter\'isticas: los pares
consulta--clave alineados punt\'uan alto y los pares ortogonales punt\'uan casi
cero.

\paragraph{La geometr\'ia de los espacios Q, K, V.}
La compatibilidad de todos los pares consulta--clave se re\'une en la
\emph{matriz de similitud}
\[
  S = Q K^\top, \qquad S_{ij} = \inner{q_i}{k_j} \in \R^{n\times n}.
\]
Esta matriz codifica todas las relaciones geom\'etricas entre consultas y
claves; cuando $Q = K$ (auto-atenci\'on pura) coincide con la matriz de Gram de
las claves. Admite una descomposici\'on en valores singulares
$S = U\Sigma V^\top$ cuyos factores capturan las direcciones principales de
consultas y claves, su norma de Frobenius satisface
$\norm{S}_F^2 = \operatorname{tr}(QK^\top K Q^\top)$ y, entrada a entrada,
$S_{ij} = \norm{q_i}\norm{k_j}\cos\theta_{ij}$.

\begin{theorem}[Proyecci\'on geom\'etrica en atenci\'on]
\label{thm:apps:geom-projection}
El mecanismo de atenci\'on puede interpretarse como una proyecci\'on geom\'etrica
del espacio de valores sobre el subespacio generado por las consultas, mediada
por la m\'etrica inducida por las claves. Para cada consulta $q_i$,
\[
  y_i = \sum_{j=1}^{n} \alpha_{ij}\, v_j = V^\top \boldsymbol{\alpha}_i,
  \qquad
  \boldsymbol{\alpha}_i = \softmax\!\bigl(S_{i,:}/\sqrt{d_k}\bigr),
\]
donde $\boldsymbol{\alpha}_i$ es una distribuci\'on de probabilidad sobre el
\'indice de valores.
\end{theorem}

\begin{proof}
Los pesos de atenci\'on $\alpha_{ij}$ forman una combinaci\'on convexa
(probabil\'istica) sobre los valores, determinada por la proximidad geom\'etrica
de $q_i$ a cada clave $k_j$, medida por el producto interno escalado y
normalizado por softmax. La proyecci\'on resultante preserva la estructura del
espacio de valores mientras enfatiza las direcciones alineadas con la consulta
seg\'un la m\'etrica definida por las claves.
\end{proof}

La \Cref{fig:apps:qkv} ilustra esta imagen en dos etapas: las similitudes se
leen como \'angulos en el espacio de consultas y claves, y la salida es la
correspondiente combinaci\'on convexa de los vectores de valores.

\begin{figure}[h]
\centering
\begin{tikzpicture}[>=Stealth, font=\small, scale=0.82, every node/.style={transform shape}]
% ============================================================
% PANEL IZQUIERDO: espacio de consultas y claves
% ============================================================
\begin{scope}[shift={(0,0)}]
  \node[font=\bfseries, align=center, text=black] at (3.2, 6.8)
    {Espacio de consultas y claves};
  \draw[->, gdlgray!60, line width=0.8pt] (-0.3, 0) -- (6.8, 0)
    node[right, font=\footnotesize, text=gdlgray!70!black] {$e_1$};
  \draw[->, gdlgray!60, line width=0.8pt] (0, -0.3) -- (0, 6.5)
    node[above, font=\footnotesize, text=gdlgray!70!black] {$e_2$};
  \foreach \x in {1,...,6} {\draw[gdlgray!15] (\x, -0.1) -- (\x, 6.3);}
  \foreach \y in {1,...,6} {\draw[gdlgray!15] (-0.1, \y) -- (6.6, \y);}
  \node[font=\footnotesize, text=gdlgray!70!black, anchor=north east] at (-0.1, -0.1) {$0$};
  \draw[gdlpurple!60, line width=0.8pt]
    ({2.2*cos(42.0)}, {2.2*sin(42.0)}) arc[start angle=42.0, end angle=51.1, radius=2.2];
  \node[font=\scriptsize, text=gdlpurple!70!black, anchor=south west]
    at ({2.55*cos(46.5)}, {2.55*sin(46.5)}) {$\theta_{i1}$};
  \draw[gdlpurple!60, line width=0.8pt]
    ({1.8*cos(23.3)}, {1.8*sin(23.3)}) arc[start angle=23.3, end angle=42.0, radius=1.8];
  \node[font=\scriptsize, text=gdlpurple!70!black, anchor=north west]
    at ({2.15*cos(32.5)}, {2.15*sin(32.5)}) {$\theta_{i2}$};
  \draw[gdlpurple!60, line width=0.8pt]
    ({1.5*cos(42.0)}, {1.5*sin(42.0)}) arc[start angle=42.0, end angle=78.3, radius=1.5];
  \node[font=\scriptsize, text=gdlpurple!70!black, anchor=south]
    at ({1.9*cos(60)}, {1.9*sin(60)}) {$\theta_{i3}$};
  \draw[-{Stealth[length=5pt,width=4pt]}, gdlblue!80, line width=1.8pt] (0,0) -- (5.0, 4.5);
  \node[font=\small, text=gdlblue!70!black, anchor=south west] at (5.1, 4.5) {$\mathbf{q}_i$};
  \fill[gdlblue!70] (5.0, 4.5) circle (3pt);
  \draw[-{Stealth[length=5pt,width=4pt]}, gdlred!80, line width=1.8pt] (0,0) -- (4.2, 5.2);
  \node[font=\small, text=gdlred!70!black, anchor=south] at (4.2, 5.35) {$\mathbf{k}_1$};
  \fill[gdlred!70] (4.2, 5.2) circle (3pt);
  \draw[-{Stealth[length=5pt,width=4pt]}, gdlred!80, line width=1.8pt] (0,0) -- (5.8, 2.5);
  \node[font=\small, text=gdlred!70!black, anchor=west] at (5.9, 2.5) {$\mathbf{k}_2$};
  \fill[gdlred!70] (5.8, 2.5) circle (3pt);
  \draw[-{Stealth[length=5pt,width=4pt]}, gdlred!80, line width=1.8pt] (0,0) -- (1.2, 5.8);
  \node[font=\small, text=gdlred!70!black, anchor=south] at (1.2, 5.9) {$\mathbf{k}_3$};
  \fill[gdlred!70] (1.2, 5.8) circle (3pt);
  \coordinate (projfoot) at (4.17, 5.17);
  \draw[gdlpurple!40, line width=0.6pt, densely dotted] (5.0, 4.5) -- (projfoot);
  \fill[gdlpurple!50] (projfoot) circle (2pt);
  \node[font=\tiny, text=gdlpurple!70!black, fill=white, fill opacity=0.85,
        text opacity=1, inner sep=1.5pt, rounded corners=1pt, anchor=south west]
    at (4.5, 5.35) {alta similitud};
  \node[font=\tiny, text=gdlpurple!70!black, fill=white, fill opacity=0.85,
        text opacity=1, inner sep=1.5pt, rounded corners=1pt, anchor=south west]
    at (6.0, 2.55) {media};
  \node[font=\tiny, text=gdlpurple!70!black, fill=white, fill opacity=0.85,
        text opacity=1, inner sep=1.5pt, rounded corners=1pt, anchor=south east]
    at (0.9, 5.95) {baja};
  \node[font=\scriptsize, text=gdlpurple!70!black, align=center, anchor=north west]
    at (0.3, -0.8) {$\alpha_{ij} = \operatorname{softmax}_j\!\left(\dfrac{\mathbf{q}_i \cdot \mathbf{k}_j}{\sqrt{d_k}}\right)$};
\end{scope}
% ============================================================
% FLECHA entre paneles
% ============================================================
\draw[->, gdlorange!70, line width=2pt] (7.8, 3.2) -- (9.5, 3.2);
\node[font=\footnotesize, align=center, text=gdlorange!70!black] at (8.65, 4.0)
  {Pesos de\\atenci\'on $\alpha_{ij}$};
% ============================================================
% PANEL DERECHO: agregaci\'on ponderada de valores
% ============================================================
\begin{scope}[shift={(10.0, 0)}]
  \node[font=\bfseries, align=center, text=black] at (3.5, 6.8)
    {Agregaci\'on ponderada de valores};
  \draw[->, gdlgray!60, line width=0.8pt] (-0.3, 0) -- (7.3, 0)
    node[right, font=\footnotesize, text=gdlgray!70!black] {$e_1$};
  \draw[->, gdlgray!60, line width=0.8pt] (0, -0.3) -- (0, 6.5)
    node[above, font=\footnotesize, text=gdlgray!70!black] {$e_2$};
  \foreach \x in {1,...,7} {\draw[gdlgray!15] (\x, -0.1) -- (\x, 6.3);}
  \foreach \y in {1,...,6} {\draw[gdlgray!15] (-0.1, \y) -- (7.1, \y);}
  \node[font=\footnotesize, text=gdlgray!70!black, anchor=north east] at (-0.1, -0.1) {$0$};
  \draw[-{Stealth[length=4pt,width=3pt]}, gdlgreen!35, line width=0.8pt] (0,0) -- (3.0, 5.5);
  \node[font=\small, text=gdlgreen!70!black, anchor=south] at (3.0, 5.6) {$\mathbf{v}_1$};
  \fill[gdlgreen!60] (3.0, 5.5) circle (2.5pt);
  \draw[-{Stealth[length=4pt,width=3pt]}, gdlgreen!35, line width=0.8pt] (0,0) -- (5.5, 2.0);
  \node[font=\small, text=gdlgreen!70!black, anchor=north west] at (5.6, 1.9) {$\mathbf{v}_2$};
  \fill[gdlgreen!60] (5.5, 2.0) circle (2.5pt);
  \draw[-{Stealth[length=3pt,width=2.5pt]}, gdlgreen!25, line width=0.6pt] (0,0) -- (1.5, 4.0);
  \node[font=\small, text=gdlgreen!70!black, anchor=east] at (1.3, 4.3) {$\mathbf{v}_3$};
  \fill[gdlgreen!50] (1.5, 4.0) circle (2pt);
  \draw[-{Stealth[length=5pt,width=4pt]}, gdlgreen!75, line width=2.2pt] (0,0) -- (1.95, 3.575);
  \node[font=\scriptsize, text=gdlgreen!70!black, anchor=east,
        fill=white, fill opacity=0.85, text opacity=1, inner sep=1pt, rounded corners=1pt]
    at (1.5, 2.5) {$\alpha_{i1}\mathbf{v}_1$};
  \draw[-{Stealth[length=4pt,width=3pt]}, gdlgreen!65, line width=1.4pt] (0,0) -- (1.375, 0.5);
  \node[font=\scriptsize, text=gdlgreen!70!black, anchor=north west,
        fill=white, fill opacity=0.85, text opacity=1, inner sep=1pt, rounded corners=1pt]
    at (1.45, 0.55) {$\alpha_{i2}\mathbf{v}_2$};
  \draw[-{Stealth[length=3pt,width=2pt]}, gdlgreen!50, line width=0.6pt] (0,0) -- (0.15, 0.4);
  \draw[gdlorange!60, line width=0.7pt, dashed] (1.95, 3.575) -- (3.325, 4.075);
  \draw[gdlorange!60, line width=0.7pt, dashed] (1.375, 0.5) -- (3.325, 4.075);
  \draw[gdlorange!50, line width=0.5pt, densely dotted] (3.325, 4.075) -- (3.475, 4.475);
  \draw[-{Stealth[length=6pt,width=5pt]}, gdlorange!80, line width=2.4pt] (0,0) -- (3.475, 4.475);
  \node[font=\small\bfseries, text=gdlorange!70!black, anchor=south west] at (3.55, 4.5) {$\mathbf{y}_i$};
  \fill[gdlorange!80] (3.475, 4.475) circle (3.5pt);
  \node[font=\scriptsize, text=gdlpurple!70!black, fill=white, fill opacity=0.85,
        text opacity=1, inner sep=2pt, rounded corners=1pt, anchor=west]
    at (3.3, 5.5) {$\alpha_{i1} = 0{,}65$};
  \node[font=\scriptsize, text=gdlpurple!70!black, fill=white, fill opacity=0.85,
        text opacity=1, inner sep=2pt, rounded corners=1pt, anchor=west]
    at (5.8, 2.5) {$\alpha_{i2} = 0{,}25$};
  \node[font=\scriptsize, text=gdlpurple!70!black, fill=white, fill opacity=0.85,
        text opacity=1, inner sep=2pt, rounded corners=1pt, anchor=east]
    at (1.0, 4.6) {$\alpha_{i3} = 0{,}10$};
  \node[font=\scriptsize, text=gdlorange!70!black, align=center, anchor=north west]
    at (0.3, -0.8) {$\mathbf{y}_i = \displaystyle\sum_{j} \alpha_{ij}\, \mathbf{v}_j$};
\end{scope}
\end{tikzpicture}
\caption[Interpretaci\'on geom\'etrica de la atenci\'on]{Interpretaci\'on
geom\'etrica del mecanismo de atenci\'on. \textbf{Izquierda}: en el espacio de
consultas y claves, la similitud entre $\mathbf{q}_i$ y cada $\mathbf{k}_j$ se
mide por el \'angulo $\theta_{ij}$ (producto interno escalado).
\textbf{Derecha}: los valores $\mathbf{v}_j$ se ponderan por los coeficientes de
atenci\'on $\alpha_{ij}$, y la salida $\mathbf{y}_i$ es su combinaci\'on convexa.}
\label{fig:apps:qkv}
\end{figure}

\subsection{El mecanismo de atenci\'on b\'asico}
\label{subsec:apps:attn-basic}

La atenci\'on de producto punto escalado se define formalmente como
\[
  \Attn(Q, K, V) = \softmax\!\left(\frac{QK^\top}{\sqrt{d_k}}\right) V .
\]
La matriz de \emph{scores} $S = QK^\top \in \R^{n\times n}$ codifica la
compatibilidad de cada par consulta--clave, y el factor de escala
$1/\sqrt{d_k}$ impide que los productos punto crezcan en exceso con la
dimensi\'on y saturen el softmax. De forma equivalente, para cada elemento
\[
  y_i = \sum_{j=1}^{n} \alpha_{ij}\, v_j,
  \qquad \sum_{j=1}^{n}\alpha_{ij} = 1, \qquad \alpha_{ij}\ge 0,
\]
de modo que la salida es una interpolaci\'on convexa en el espacio de valores,
con pesos que forman una distribuci\'on de probabilidad sobre las entradas.

\begin{definition}[Funci\'on de atenci\'on generalizada]
\label{def:apps:generalized-attn}
Una \emph{funci\'on de atenci\'on generalizada} calcula, para cada consulta,
\[
  y_i = \sum_{j=1}^{n}\alpha_{ij}\, v_j,
  \qquad
  \alpha_{ij} = \frac{\exp(e_{ij})}{\sum_{k=1}^n \exp(e_{ik})},
  \qquad
  e_{ij} = a(q_i, k_j),
\]
donde $a:\mathcal{X}\times\mathcal{Y}\to\R$ es una funci\'on de compatibilidad.
El producto punto escalado es la elecci\'on $a(q,k) = \inner{q}{k}/\sqrt{d_k}$,
pero cualquier escalar $G$-invariante del par $(q_i, k_j)$ es admisible ---la
libertad que explotan las variantes geom\'etricas siguientes.
\end{definition}

La \emph{auto-atenci\'on} es el caso especial $Q = K = V = X$ en el que
consultas, claves y valores provienen de la misma entrada, lo que permite que
cada elemento atienda a todos los dem\'as y capture dependencias globales en una
sola operaci\'on. Esta construcci\'on es adem\'as m\'aximamente expresiva entre
las funciones permutaci\'on-equivariantes.

\begin{theorem}[Universalidad de la auto-atenci\'on]
\label{thm:apps:self-attn-universal}
Sea $\mathcal{K}\subseteq\R^{n\times d}$ compacto. Toda funci\'on continua
permutaci\'on-equivariante $f:\mathcal{K}\to\R^{n\times d}$ (es decir,
$f(PX) = P f(X)$ para toda matriz de permutaci\'on $P$) puede aproximarse
uniformemente en $\mathcal{K}$ con precisi\'on arbitraria. La universalidad plena
requiere el bloque \emph{transformer} completo ---auto-atenci\'on junto con capas
\emph{feed-forward} y conexiones residuales---, no capas de auto-atenci\'on
aisladas; v\'ease \citet{yun2020transformers} para el enunciado riguroso.
\end{theorem}

\begin{proof}[Esquema]
Por el teorema de representaci\'on para funciones sim\'etricas~\citep{zaheer2017deepsets},
toda $f$ continua y permutaci\'on-equivariante puede escribirse como
$f(X)_i = \phi\bigl(x_i, \sum_j \rho(x_j)\bigr)$ para funciones continuas
$\phi,\rho$ apropiadas. Por el teorema de aproximaci\'on universal se aproxima
esto, para todo $\epsilon>0$, mediante una suma de interacciones por pares
$\sum_j \softmax\bigl(g_1(x_i)^\top g_2(x_j)\bigr)\, h(g_3(x_j))$, que se realiza
exactamente con una capa de auto-atenci\'on cuyas matrices de proyecci\'on
implementan $g_1,g_2,g_3$, absorbiendo el factor de escala en $g_1,g_2$. La
composici\'on de funciones universalmente aproximables da el resultado.
\end{proof}

\subsection{Atenci\'on multicabeza}
\label{subsec:apps:multihead}

La atenci\'on multicabeza extiende el mecanismo b\'asico procesando varios
subespacios de representaci\'on en paralelo,
\[
  \operatorname{MultiHead}(Q,K,V) = \Concat(\mathrm{head}_1,\dots,\mathrm{head}_h)\,W_O,
  \qquad
  \mathrm{head}_i = \Attn(QW_Q^i, KW_K^i, VW_V^i),
\]
con proyecciones $W_Q^i, W_K^i \in \R^{d\times d_k}$ y
$W_V^i \in \R^{d\times d_v}$ para $i = 1,\dots,h$ y una proyecci\'on de salida
$W_O \in \R^{hd_v\times d}$. La descomposici\'on en subespacios paralelos permite
que cada cabeza se especialice en un tipo distinto de relaci\'on, aumenta la
capacidad del modelo para atender a varias posiciones y subespacios a la vez y
mejora la estabilidad del entrenamiento gracias a la redundancia parcial
(\Cref{fig:apps:multihead}).

\begin{figure}[h]
\centering
\begin{tikzpicture}[
    scale=0.78, every node/.style={transform shape},
    arrow/.style={-{Stealth[length=5pt]}, thick, draw=gdlgray!60},
    line/.style={thick, draw=gdlgray!60},
    inputbox/.style={rectangle, rounded corners=3pt, draw=gdlblue!60, fill=gdlblue!20,
        thick, minimum width=2.4cm, minimum height=0.7cm, font=\small,
        text=gdlblue!70!black, align=center},
    projbox/.style={rectangle, rounded corners=3pt, draw=gdlgreen!60, fill=gdlgreen!20,
        thick, minimum width=1.1cm, minimum height=0.55cm, font=\footnotesize,
        text=gdlgreen!70!black, align=center},
    attnbox/.style={rectangle, rounded corners=3pt, draw=gdlorange!60, fill=gdlorange!20,
        thick, minimum width=2.8cm, minimum height=0.6cm, font=\footnotesize,
        text=gdlorange!70!black, align=center},
    concatbox/.style={rectangle, rounded corners=3pt, draw=gdlpurple!60, fill=gdlpurple!20,
        thick, minimum width=3.2cm, minimum height=0.7cm, font=\small,
        text=gdlpurple!70!black, align=center},
    outputbox/.style={rectangle, rounded corners=3pt, draw=gdlblue!60, fill=gdlblue!20,
        thick, minimum width=2.4cm, minimum height=0.7cm, font=\small,
        text=gdlblue!70!black, align=center},
    headlabel/.style={font=\footnotesize\bfseries, text=gdlgray!70!black},
    dashedbox/.style={rectangle, rounded corners=5pt, draw=gdlgray!40, dashed,
        thick, inner sep=10pt, inner ysep=10pt}
]
\def\headspacing{4.0}
\def\projy{-2.8}
\def\attny{-4.2}
\def\distribY{-1.2}
\def\concaty{-6.0}
\def\outputy{-7.5}
\node[inputbox] (input) at (0, 0) {$\mathbf{X} \in \mathbb{R}^{n \times d}$};
\node[projbox] (q1) at ({-\headspacing - 0.65}, \projy) {$\mathbf{Q}_1$};
\node[projbox] (k1) at ({-\headspacing + 0.35}, \projy) {$\mathbf{K}_1$};
\node[projbox] (v1) at ({-\headspacing + 1.35}, \projy) {$\mathbf{V}_1$};
\node[attnbox] (attn1) at ({-\headspacing + 0.35}, \attny)
    {$\mathrm{Attn}(\mathbf{Q}_1, \mathbf{K}_1, \mathbf{V}_1)$};
\node[projbox] (q2) at ({-0.65}, \projy) {$\mathbf{Q}_2$};
\node[projbox] (k2) at (0.35, \projy) {$\mathbf{K}_2$};
\node[projbox] (v2) at (1.35, \projy) {$\mathbf{V}_2$};
\node[attnbox] (attn2) at (0.35, \attny)
    {$\mathrm{Attn}(\mathbf{Q}_2, \mathbf{K}_2, \mathbf{V}_2)$};
\node[projbox] (q3) at ({\headspacing - 0.65}, \projy) {$\mathbf{Q}_3$};
\node[projbox] (k3) at ({\headspacing + 0.35}, \projy) {$\mathbf{K}_3$};
\node[projbox] (v3) at ({\headspacing + 1.35}, \projy) {$\mathbf{V}_3$};
\node[attnbox] (attn3) at ({\headspacing + 0.35}, \attny)
    {$\mathrm{Attn}(\mathbf{Q}_3, \mathbf{K}_3, \mathbf{V}_3)$};
\foreach \h in {1,2,3} {
    \draw[arrow] (q\h.south) -- (attn\h.north -| q\h.south);
    \draw[arrow] (k\h.south) -- (attn\h.north -| k\h.south);
    \draw[arrow] (v\h.south) -- (attn\h.north -| v\h.south);
}
\node[dashedbox, fit=(q1)(v1)(attn1)] (hbox1) {};
\node[dashedbox, fit=(q2)(v2)(attn2)] (hbox2) {};
\node[dashedbox, fit=(q3)(v3)(attn3)] (hbox3) {};
\node[headlabel, anchor=south] at (hbox1.north) {Cabeza 1};
\node[headlabel, anchor=south] at (hbox2.north) {Cabeza 2};
\node[headlabel, anchor=south] at (hbox3.north) {Cabeza 3};
\draw[line] (input.south) -- (0, \distribY);
\draw[line] ({-\headspacing + 0.35}, \distribY) -- ({\headspacing + 0.35}, \distribY);
\pgfmathsetmacro{\arrowEndY}{\projy + 1.45}
\draw[arrow] ({-\headspacing + 0.35}, \distribY) -- ({-\headspacing + 0.35}, \arrowEndY);
\draw[arrow] (0.35, \distribY) -- (0.35, \arrowEndY);
\draw[arrow] ({\headspacing + 0.35}, \distribY) -- ({\headspacing + 0.35}, \arrowEndY);
\node[concatbox] (concat) at (0.35, \concaty)
    {$\mathrm{Concat}(\mathrm{head}_1, \ldots, \mathrm{head}_h)$};
\pgfmathsetmacro{\convY}{\concaty + 0.7}
\draw[line] (attn1.south) -- (attn1.south |- 0, \convY);
\draw[arrow] (attn1.south |- 0, \convY) -| ($(concat.north) + (-0.5, 0)$);
\draw[arrow] (attn2.south) -- (concat.north);
\draw[line] (attn3.south) -- (attn3.south |- 0, \convY);
\draw[arrow] (attn3.south |- 0, \convY) -| ($(concat.north) + (0.5, 0)$);
\node[outputbox] (output) at (0.35, \outputy) {$\mathbf{W}^O \cdot \mathrm{Concat}$};
\draw[arrow] (concat.south) -- (output.north);
\node[font=\small, text=gdlblue!70!black, below=4pt of output]
    {Salida $\in \mathbb{R}^{n \times d}$};
\node[font=\scriptsize, text=gdlgray!70!black, align=left, anchor=north west]
    at ($(hbox3.east) + (0.5, 0.4)$)
    {$\mathbf{W}_i^Q, \mathbf{W}_i^K \in \mathbb{R}^{d \times d_k}$\\[2pt]
     $\mathbf{W}_i^V \in \mathbb{R}^{d \times d_v}$\\[2pt]
     $\mathbf{W}^O \in \mathbb{R}^{hd_v \times d}$\\[2pt]
     $d_k = d_v = d/h$};
\end{tikzpicture}
\caption[Atenci\'on multicabeza]{Atenci\'on multicabeza. La entrada se proyecta
en $h$ subespacios paralelos, cada uno con sus propias matrices
$\mathbf{W}^Q_i$, $\mathbf{W}^K_i$, $\mathbf{W}^V_i$. Cada cabeza computa la
atenci\'on de forma independiente y los resultados se concatenan y se proyectan
mediante $\mathbf{W}^O$.}
\label{fig:apps:multihead}
\end{figure}

\begin{proposition}[Capacidad representacional de los \emph{scores} multicabeza pre-softmax]
\label{prop:apps:multihead-capacity}
La atenci\'on multicabeza con $h$ cabezas puede representar cualquier patr\'on de
\emph{scores} pre-softmax que admita descomposici\'on como suma de $h$ t\'erminos
de rango bajo. Formalmente, si $S^\star = \sum_{i=1}^h U_i\Sigma_i V_i^\top$ con
cada $U_i\Sigma_i V_i^\top$ de rango $r_i\le d_k$, entonces existe una
configuraci\'on con $h$ cabezas cuya matriz de \emph{scores} pre-softmax es
exactamente $S^\star$.
\end{proposition}

\begin{proof}
Tomamos $W_Q^i = U_i\Sigma_i^{1/2}$, $W_K^i = V_i\Sigma_i^{1/2}$ y $W_V^i = I$
para cada t\'ermino. Entonces la cabeza $i$ tiene \emph{scores} pre-softmax
$Q_i K_i^\top = X U_i\Sigma_i V_i^\top X^\top$, y la suma sobre las cabezas
reproduce $S^\star$.
\end{proof}

\begin{remark}
\label{rem:apps:multihead-nonlinearity}
Esta construcci\'on opera a nivel de \emph{scores} pre-softmax. En la pr\'actica
cada cabeza aplica softmax antes de agregar con los valores, lo que obliga a que
los pesos de atenci\'on efectivos sean no negativos y sumen uno por fila. Esta no
linealidad impide que el mecanismo multicabeza represente exactamente funciones
de atenci\'on arbitrarias a nivel de la salida, pero el resultado de capacidad
explica por qu\'e unas pocas cabezas ya capturan una estructura multirelacional
rica. El coste en tiempo y espacio de la capa es $O(n^2 d + n d h d_k)$ y
$O(n^2 + ndh)$ respectivamente; el t\'ermino $n^2$ domina para secuencias largas
y motiva las aproximaciones dispersas y lineales.
\end{remark}

\subsection{La atenci\'on como operador integral}
\label{subsec:apps:integral-operator}

La auto-atenci\'on puede formalizarse rigurosamente como un operador integral que
act\'ua sobre funciones definidas en una variedad, una perspectiva que la conecta
con la teor\'ia cl\'asica de operadores compactos y expone su comportamiento
espectral. Sea $\mathcal{M}$ una variedad riemanniana compacta dotada de una
medida $\mu$. Para una funci\'on $f:\mathcal{M}\to\R^d$ el operador de atenci\'on
es
\[
  \mathcal{T}f(x) = \int_{\mathcal{M}} \kappa(x,y)\, f(y)\, d\mu(y),
  \qquad
  \kappa(x,y) = \sigma\!\left(\frac{\inner{Q(x)}{K(y)}}{\sqrt{d_k}}\right) V(y)^\top,
\]
con $Q,K,V$ aplicando $\mathcal{M}$ en los espacios de consultas, claves y
valores y $\sigma$ un softmax generalizado. El \emph{transformer} discreto es la
cuadratura emp\'irica de esta integral sobre $n$ puntos de muestreo.

\begin{proposition}[Compacidad del operador de atenci\'on]
\label{prop:apps:attn-compact}
Si el n\'ucleo $\kappa(x,y)$ es continuo en $\mathcal{M}\times\mathcal{M}$ y
$\mathcal{M}$ es compacta, entonces
$\mathcal{T}: L^2(\mathcal{M})\to L^2(\mathcal{M})$ es compacto.
\end{proposition}

\begin{proof}[Esquema]
Para $f$ en la bola unidad de $L^2(\mathcal{M})$, la imagen $\mathcal{T}f$ est\'a
uniformemente acotada (pues $\kappa$ alcanza su m\'aximo en el compacto
$\mathcal{M}\times\mathcal{M}$) y es equicontinua (pues $\kappa$ es
uniformemente continua all\'i). Por Arzel\`a--Ascoli la imagen de la bola unidad
es relativamente compacta en $C(\mathcal{M})$ y, por tanto, en
$L^2(\mathcal{M})$, de modo que $\mathcal{T}$ es compacto.
\end{proof}

Cuando el n\'ucleo es un n\'ucleo de Mercer ---sim\'etrico y definido positivo,
como en la variante gaussiana
$\kappa_{\text{gauss}}(x,y) = \exp(-\norm{Q(x)-K(y)}^2/2\sigma^2)$--- el teorema
espectral para operadores compactos autoadjuntos proporciona una descomposici\'on
$\kappa(x,y) = \sum_{i} \lambda_i\,\phi_i(x)\phi_i(y)$ con autovalores no
negativos $\lambda_1\ge\lambda_2\ge\cdots\ge 0$ y autofunciones ortonormales
$\phi_i$, de modo que
\[
  \mathcal{T}f = \sum_{i=1}^{\infty}\lambda_i \inner{f}{\phi_i}_{L^2}\phi_i .
\]
Las autofunciones son los ``modos de atenci\'on'' fundamentales y los autovalores
los ordenan por importancia (\Cref{fig:apps:spectrum}). Como los autovalores
suelen decaer r\'apidamente, la atenci\'on act\'ua como un filtro pasa-bajas
natural: la truncaci\'on de rango $r$, $\mathcal{T}_r$, satisface la cota de error
$\norm{\mathcal{T}f - \mathcal{T}_r f}_{L^2}\le \lambda_{r+1}\norm{f}_{L^2}$, de
modo que un rango efectivo finito ya captura el operador. Cuando el n\'ucleo es
invariante a traslaciones, esta descomposici\'on espectral se reduce al
an\'alisis cl\'asico de Fourier del \Cref{ch:grids}, con autofunciones
exponenciales complejas y autovalores la transformada de Fourier del n\'ucleo
---un puente directo entre el operador de atenci\'on y la convoluci\'on.

\begin{figure}[h]
\centering
\begin{tikzpicture}[>=Stealth, font=\small, scale=0.78, every node/.style={transform shape}]
% PANEL IZQUIERDO: espectro de autovalores
\begin{scope}[shift={(0,0)}]
  \node[font=\bfseries, text=black] at (4.0, 6.2) {Espectro del operador de atenci\'on};
  \fill[gdlyellow!10] (0.27, 0) rectangle (3.65, 5.2);
  \draw[->, gdlgray!70, line width=0.8pt] (-0.3, 0) -- (8.5, 0)
    node[right, font=\footnotesize, text=gdlgray!70!black] {$i$};
  \draw[->, gdlgray!70, line width=0.8pt] (0, -0.3) -- (0, 5.8)
    node[above, font=\footnotesize, text=gdlgray!70!black] {$\lambda_i$};
  \foreach \y/\lab in {1/0{,}2, 2/0{,}4, 3/0{,}6, 4/0{,}8, 5/1{,}0} {
    \draw[gdlgray!50] (-0.1, \y) -- (0.1, \y);
    \node[font=\tiny, text=gdlgray!70!black, anchor=east] at (-0.15, \y) {$\lab$};}
  \def\barwidth{0.45}
  \fill[gdlblue!60] (0.5-\barwidth/2, 0) rectangle (0.5+\barwidth/2, 5.0);
  \node[font=\tiny, text=gdlgray!70!black] at (0.5, -0.3) {$1$};
  \fill[gdlblue!55] (1.2-\barwidth/2, 0) rectangle (1.2+\barwidth/2, 2.5);
  \node[font=\tiny, text=gdlgray!70!black] at (1.2, -0.3) {$2$};
  \fill[gdlblue!50] (1.9-\barwidth/2, 0) rectangle (1.9+\barwidth/2, 1.48);
  \node[font=\tiny, text=gdlgray!70!black] at (1.9, -0.3) {$3$};
  \fill[gdlblue!45] (2.6-\barwidth/2, 0) rectangle (2.6+\barwidth/2, 1.0);
  \node[font=\tiny, text=gdlgray!70!black] at (2.6, -0.3) {$4$};
  \fill[gdlblue!40] (3.3-\barwidth/2, 0) rectangle (3.3+\barwidth/2, 0.72);
  \node[font=\tiny, text=gdlgray!70!black] at (3.3, -0.3) {$5$};
  \fill[gdlblue!35] (4.0-\barwidth/2, 0) rectangle (4.0+\barwidth/2, 0.55);
  \node[font=\tiny, text=gdlgray!70!black] at (4.0, -0.3) {$6$};
  \fill[gdlblue!30] (4.7-\barwidth/2, 0) rectangle (4.7+\barwidth/2, 0.43);
  \node[font=\tiny, text=gdlgray!70!black] at (4.7, -0.3) {$7$};
  \fill[gdlblue!25] (5.4-\barwidth/2, 0) rectangle (5.4+\barwidth/2, 0.35);
  \node[font=\tiny, text=gdlgray!70!black] at (5.4, -0.3) {$8$};
  \fill[gdlblue!20] (6.1-\barwidth/2, 0) rectangle (6.1+\barwidth/2, 0.28);
  \node[font=\tiny, text=gdlgray!70!black] at (6.1, -0.3) {$9$};
  \fill[gdlblue!18] (6.8-\barwidth/2, 0) rectangle (6.8+\barwidth/2, 0.24);
  \node[font=\tiny, text=gdlgray!70!black] at (6.8, -0.3) {$10$};
  \node[font=\small, text=gdlgray!70!black] at (7.6, 0.15) {$\cdots$};
  \node[font=\tiny, text=gdlgray!70!black] at (7.6, -0.3) {$n$};
  \draw[gdlred!70, line width=1.2pt, dashed, smooth]
    plot[domain=0.5:8.0, samples=50] (\x, {5.0/(\x)^1.5});
  \node[font=\scriptsize, text=gdlred!70!black, anchor=west, fill=white,
        fill opacity=0.85, text opacity=1, inner sep=2pt, rounded corners=1pt]
    at (5.5, 1.8) {$\lambda_i \sim i^{-\alpha}$};
  \draw[gdlorange!70, line width=1.0pt, densely dash dot] (3.65, -0.15) -- (3.65, 5.5);
  \node[font=\scriptsize, text=gdlorange!70!black, anchor=south, fill=white,
        fill opacity=0.85, text opacity=1, inner sep=2pt, rounded corners=1pt]
    at (3.65, 5.5) {rango efectivo $r$};
  \node[font=\scriptsize, text=gdlblue!70!black, align=center, anchor=north]
    at (4.0, -0.9) {$\kappa(x,y) = \displaystyle\sum_{i=1}^{\infty} \lambda_i \,\phi_i(x)\,\phi_i(y)$};
\end{scope}
% PANEL DERECHO: autofunciones
\begin{scope}[shift={(10.5, 0)}]
  \node[font=\bfseries, text=black] at (2.5, 6.2) {Autofunciones $\phi_i$};
  \node[font=\scriptsize, text=gdlblue!70!black, anchor=east] at (-0.3, 4.8) {$\phi_1$};
  \draw[->, gdlgray!50, line width=0.4pt] (0, 4.8) -- (5.2, 4.8);
  \draw[gdlgreen!70, line width=1.4pt, smooth]
    plot[domain=0:5, samples=40] (\x, {4.8 + 0.4*cos(18*\x)});
  \node[font=\tiny, text=gdlpurple!70!black, anchor=west] at (5.3, 4.8) {$\lambda_1$};
  \fill[gdlpurple!60] (5.9, 4.72) rectangle (6.9, 4.88);
  \node[font=\scriptsize, text=gdlblue!70!black, anchor=east] at (-0.3, 3.3) {$\phi_2$};
  \draw[->, gdlgray!50, line width=0.4pt] (0, 3.3) -- (5.2, 3.3);
  \draw[gdlgreen!70, line width=1.4pt, smooth]
    plot[domain=0:5, samples=60] (\x, {3.3 + 0.45*cos(72*\x)});
  \node[font=\tiny, text=gdlpurple!70!black, anchor=west] at (5.3, 3.3) {$\lambda_2$};
  \fill[gdlpurple!55] (5.9, 3.22) rectangle (6.4, 3.38);
  \node[font=\scriptsize, text=gdlblue!70!black, anchor=east] at (-0.3, 1.8) {$\phi_3$};
  \draw[->, gdlgray!50, line width=0.4pt] (0, 1.8) -- (5.2, 1.8);
  \draw[gdlgreen!70, line width=1.4pt, smooth]
    plot[domain=0:5, samples=80] (\x, {1.8 + 0.45*sin(144*\x)});
  \node[font=\tiny, text=gdlpurple!70!black, anchor=west] at (5.3, 1.8) {$\lambda_3$};
  \fill[gdlpurple!50] (5.9, 1.72) rectangle (6.2, 1.88);
  \node[font=\scriptsize, text=gdlblue!70!black, anchor=east] at (-0.3, 0.3) {$\phi_4$};
  \draw[->, gdlgray!50, line width=0.4pt] (0, 0.3) -- (5.2, 0.3);
  \draw[gdlgreen!70, line width=1.4pt, smooth]
    plot[domain=0:5, samples=100] (\x, {0.3 + 0.4*cos(216*\x)});
  \node[font=\tiny, text=gdlpurple!70!black, anchor=west] at (5.3, 0.3) {$\lambda_4$};
  \fill[gdlpurple!45] (5.9, 0.22) rectangle (6.1, 0.38);
  \node[font=\small, text=gdlgray!60] at (2.5, -0.4) {$\vdots$};
  \node[font=\scriptsize, text=gdlgreen!60!black, align=center, anchor=north]
    at (2.5, -0.8) {Frecuencia creciente $\longrightarrow$};
  \draw[gdlpurple!50, line width=0.6pt, decorate,
        decoration={brace, amplitude=4pt, mirror}] (7.1, 0.15) -- (7.1, 4.95);
  \node[font=\tiny, text=gdlpurple!70!black, anchor=west, align=center] at (7.4, 2.55)
    {Decaimiento\\$\lambda_i \to 0$};
\end{scope}
\end{tikzpicture}
\caption[Descomposici\'on espectral de la atenci\'on]{Descomposici\'on espectral
del operador de atenci\'on. \textbf{Izquierda}: los autovalores $\lambda_i$ del
n\'ucleo de atenci\'on decaen r\'apidamente ($\lambda_i\sim i^{-\alpha}$), lo que
indica un rango efectivo finito $r$; la regi\'on sombreada corresponde a los
modos dominantes. \textbf{Derecha}: las autofunciones $\phi_i$ asociadas muestran
frecuencia creciente, an\'alogamente a una base de Fourier; las barras p\'urpura
indican la magnitud del autovalor correspondiente.}
\label{fig:apps:spectrum}
\end{figure}

\subsection{La atenci\'on como paso de mensajes adaptativo}
\label{subsec:apps:attn-mpnn}

Le\'ida geom\'etricamente, la atenci\'on de producto punto escalado es un paso de
mensajes (\Cref{ch:graphs}) en el que los pesos de las aristas no est\'an
fijados por el dominio, sino que se \emph{calculan a partir de los datos}: cada
elemento consulta a los dem\'as y decide cu\'anto escuchar a cada uno. Donde una
convoluci\'on codifica r\'igidamente los pesos del vecindario y una red de grafos
los fija por la topolog\'ia, la atenci\'on los aprende sobre la marcha. Esto
sit\'ua a la atenci\'on en el extremo adaptativo de un espectro cuyo extremo
r\'igido es la convoluci\'on equivariante ---una tensi\'on entre la estructura
impuesta y la flexibilidad aprendida que recorre todo el aprendizaje profundo
geom\'etrico y que la \Cref{sec:apps:cnn-attention} hace precisa.

\paragraph{Atenci\'on sobre grafos.}
La instancia geom\'etrica m\'as directa es la red de atenci\'on sobre grafos
(GAT)~\citep{velickovic2017graph}, que restringe la atenci\'on a los vecinos de
cada nodo y los pondera de forma adaptativa. Para un nodo $i$ con vecindario
$\mathcal{N}(i)$, calcula los coeficientes de atenci\'on a partir de las
caracter\'isticas de los nodos, los normaliza sobre el vecindario y agrega,
\[
  e_{ij} = \operatorname{LeakyReLU}\!\bigl(a^\top[Wh_i \,\|\, Wh_j]\bigr),
  \qquad
  \alpha_{ij} = \softmax_j(e_{ij}),
  \qquad
  h_i' = \sigma\!\Bigl(\textstyle\sum_{j\in\mathcal{N}(i)} \alpha_{ij}\,Wh_j\Bigr).
\]
Como los coeficientes dependen solo de las caracter\'isticas de los nodos
implicados y no de ning\'un orden, la capa es permutaci\'on-equivariante por
construcci\'on (\Cref{fig:apps:gat}). Como toda red de paso de mensajes
est\'andar, su poder para distinguir grafos est\'a acotado por el test de
Weisfeiler--Leman~\citep{xu2019how}, pero la ponderaci\'on aprendida y
anis\'otropa de los vecinos la hace en la pr\'actica notablemente m\'as expresiva
que la agregaci\'on de peso fijo.

\begin{figure}[h]
\centering
\begin{tikzpicture}[>=Stealth, font=\small]
\tikzset{
    targetnode/.style={circle, draw=gdlred!70, fill=gdlred!20,
        minimum size=28pt, inner sep=0pt, font=\small\bfseries, line width=1.6pt},
    neighnode/.style={circle, draw=gdlblue!60, fill=gdlblue!20,
        minimum size=22pt, inner sep=0pt, font=\small},
    nonneighnode/.style={circle, draw=gdlgray!40, fill=gdlgray!15,
        minimum size=20pt, inner sep=0pt, font=\small, text=gdlgray!70!black}
}
\node[targetnode] (vi) at (0, 0) {$v_i$};
\node[neighnode] (v1) at (-2.4, 1.8)  {$v_1$};
\node[neighnode] (v2) at (0.0, 2.6)   {$v_2$};
\node[neighnode] (v3) at (2.4, 1.4)   {$v_3$};
\node[neighnode] (v4) at (-2.2, -1.4) {$v_4$};
\node[nonneighnode] (v5) at (2.6, -1.6)  {$v_5$};
\node[nonneighnode] (v6) at (-3.6, 0.2)  {$v_6$};
\draw[gdlgray!40, thin] (v1) -- (v6);
\draw[gdlgray!40, thin] (v4) -- (v6);
\draw[gdlgray!40, thin] (v3) -- (v5);
\draw[gdlgray!40, thin] (v1) -- (v2);
\draw[gdlgray!40, thin] (vi) -- (v1);
\draw[gdlgray!40, thin] (vi) -- (v2);
\draw[gdlgray!40, thin] (vi) -- (v3);
\draw[gdlgray!40, thin] (vi) -- (v4);
\draw[gdlgray!40, thin] (v5) -- (v3);
\draw[-{Stealth[length=7pt,width=5pt]}, line width=3.2pt,
    gdlorange, opacity=0.90, shorten >=8pt, shorten <=7pt] (v2) -- (vi);
\node[font=\scriptsize, text=gdlorange!70!black, fill=white, inner sep=1.5pt,
    rounded corners=1pt] at ($(v2)!0.48!(vi) + (0.7, 0.1)$) {$\alpha_{2i} = 0{,}42$};
\draw[-{Stealth[length=6pt,width=4.5pt]}, line width=2.2pt,
    gdlorange, opacity=0.75, shorten >=8pt, shorten <=7pt] (v1) -- (vi);
\node[font=\scriptsize, text=gdlorange!70!black, fill=white, inner sep=1.5pt,
    rounded corners=1pt] at ($(v1)!0.48!(vi) + (-0.85, -0.15)$) {$\alpha_{1i} = 0{,}28$};
\draw[-{Stealth[length=5.5pt,width=4pt]}, line width=1.5pt,
    gdlorange, opacity=0.60, shorten >=8pt, shorten <=7pt] (v3) -- (vi);
\node[font=\scriptsize, text=gdlorange!70!black, fill=white, inner sep=1.5pt,
    rounded corners=1pt] at ($(v3)!0.48!(vi) + (0.85, -0.15)$) {$\alpha_{3i} = 0{,}19$};
\draw[-{Stealth[length=5pt,width=3.5pt]}, line width=0.9pt,
    gdlorange, opacity=0.45, shorten >=8pt, shorten <=7pt] (v4) -- (vi);
\node[font=\scriptsize, text=gdlorange!70!black, fill=white, inner sep=1.5pt,
    rounded corners=1pt] at ($(v4)!0.48!(vi) + (-0.85, 0.15)$) {$\alpha_{4i} = 0{,}11$};
\begin{scope}[shift={(-3.6, -2.8)}]
    \node[font=\scriptsize\bfseries, text=gdlgray!70!black, anchor=west]
        at (0, 0) {Grosor $\propto$ peso de atenci\'on};
    \draw[-{Stealth[length=4pt,width=3pt]}, line width=0.9pt,
        gdlorange, opacity=0.45] (0, -0.4) -- (1.2, -0.4);
    \node[font=\tiny, text=gdlgray!70!black, anchor=west] at (1.35, -0.4) {bajo};
    \draw[-{Stealth[length=5pt,width=4pt]}, line width=3.2pt,
        gdlorange, opacity=0.90] (2.4, -0.4) -- (3.6, -0.4);
    \node[font=\tiny, text=gdlgray!70!black, anchor=west] at (3.75, -0.4) {alto};
\end{scope}
\end{tikzpicture}
\caption[Mecanismo de atenci\'on sobre grafos]{El mecanismo de atenci\'on sobre
grafos. El nodo objetivo (rojo) calcula los coeficientes de atenci\'on
$\alpha_{ij}$ con cada vecino (azul); el grosor de la flecha refleja la magnitud
del peso aprendido. La red aprende a asignar una importancia distinta a cada
vecino a partir \'unicamente de las caracter\'isticas de los nodos, de modo que
la operaci\'on es invariante al orden de los vecinos.}
\label{fig:apps:gat}
\end{figure}

\paragraph{Atenci\'on equivariante en 3D.}
Cuando las entradas portan coordenadas geom\'etricas, la atenci\'on debe respetar
las simetr\'ias del espacio. El principio general es limpio: un mecanismo de
atenci\'on es $G$-equivariante precisamente cuando los \emph{scores} de
compatibilidad son $G$-invariantes y la transformaci\'on de los valores conmuta
con la acci\'on del grupo.

\begin{proposition}[Atenci\'on equivariante]
\label{prop:apps:equiv-attn}
Una capa de atenci\'on es $G$-equivariante si los \emph{scores} de atenci\'on
$\alpha_{ij}$ son invariantes bajo la acci\'on de $G$ y la aplicaci\'on de
valores conmuta con dicha acci\'on. Entonces $f(g\cdot x) = g\cdot f(x)$ para
todo $g\in G$.
\end{proposition}

El $\SE(3)$-Transformer~\citep{fuchs2020se3transformers} lo realiza para los
movimientos r\'igidos en tres dimensiones. Calcula los \emph{scores} de
atenci\'on a partir de cantidades invariantes bajo rotaciones y traslaciones
---caracter\'isticas escalares, distancias interat\'omicas $\norm{r_i - r_j}$ y
productos internos de caracter\'isticas vectoriales--- y agrega
caracter\'isticas de tipo~$\ell$ de modo que los vectores rotan con la escena
mientras los escalares permanecen fijos. El Equiformer~\citep{liao2022equiformer}
reemplaza el producto punto por una red equivariante que act\'ua sobre productos
tensoriales de representaciones irreducibles, lo que permite interacciones m\'as
ricas preservando la equivarianza $\SO(3)$, y el Point
Transformer~\citep{zhao2021point} lleva la atenci\'on vectorial a las nubes de
puntos. Estos modelos son los caballos de batalla de las aplicaciones
moleculares y f\'isicas de la \Cref{sec:apps:molecules}.

\paragraph{Atenci\'on no euclidiana y algebraica.}
Dos direcciones adicionales ampl\'ian el alcance de la atenci\'on. Las redes de
atenci\'on hiperb\'olica~\citep{gulcehre2019hyperbolic} reemplazan el producto
interno euclidiano por operaciones en el espacio hiperb\'olico, cuyo volumen
crece exponencialmente y embebe datos arb\'oreos y jer\'arquicos con mucha menos
distorsi\'on que el espacio euclidiano (desarrollado en \Cref{ch:geodesics}). Los
\emph{Geometric Algebra Transformers} (GATr)~\citep{brehmer2023geometric}
representan puntos, rectas y planos de manera uniforme como multivectores de un
\'algebra de Clifford y construyen sobre ellos una atenci\'on
$\Eucl(3)$-equivariante, dando lugar a una \'unica arquitectura que maneja
diversas primitivas geom\'etricas 3D. El hilo com\'un es que el sesgo inductivo
de la geometr\'ia ---ya sea la curvatura negativa o el \'algebra de los
movimientos r\'igidos--- se lleva al propio operador de atenci\'on.

\section{De la convoluci\'on a la atenci\'on: una visi\'on unificada}
\label{sec:apps:cnn-attention}

Los fundamentos de la atenci\'on de la secci\'on anterior y las convoluciones
geom\'etricas de los \Cref{ch:grids,ch:groups,ch:graphs} no son dos paradigmas
en competencia, sino dos puntos de un mismo continuo. Esta secci\'on precisa la
conexi\'on y cierra el arco te\'orico del libro antes de pasar a las
aplicaciones.

\begin{definition}[Paso de mensajes geom\'etrico generalizado]
\label{def:apps:geom-mp}
Sea $\mathcal{D}$ un dominio geom\'etrico (una rejilla euclidiana, un grafo, una
variedad, un espacio homog\'eneo), $G$ un grupo de simetr\'ias y
$\mathcal{F}(\mathcal{D})$ el espacio de se\~nales sobre $\mathcal{D}$. Un
operador de paso de mensajes geom\'etrico es una transformaci\'on
$\mathcal{T}^{(G)}:\mathcal{F}(\mathcal{D})\to\mathcal{F}(\mathcal{D})$ que es
$G$-equivariante ($\mathcal{T}^{(G)}(g\cdot f) = g\cdot\mathcal{T}^{(G)}(f)$),
geom\'etricamente local (su soporte est\'a determinado por la geometr\'ia de
$\mathcal{D}$) y composicional (las composiciones preservan estas propiedades).
\end{definition}

\begin{theorem}[Correspondencia CNN--atenci\'on]
\label{thm:apps:cnn-attn-correspondence}
Existe una correspondencia entre las convoluciones geom\'etricas y la atenci\'on
geom\'etrica a trav\'es del paso de mensajes geom\'etrico: una convoluci\'on
geom\'etrica es paso de mensajes con pesos \emph{fijos}, la atenci\'on
geom\'etrica es paso de mensajes con pesos \emph{adaptativos}, y las
arquitecturas h\'ibridas se sit\'uan en medio, con pesos parcialmente
estructurados. La correspondencia preserva la equivarianza, el poder expresivo y
la complejidad computacional.
\end{theorem}

\begin{proof}[Esquema]
La equivalencia se da en tres niveles. \emph{Algebraicamente}, ambas operaciones
son convoluciones generalizadas en el \'algebra de grupo $\C[G]$: para una
convoluci\'on el n\'ucleo es fijo, $k_g$; para la atenci\'on es din\'amico,
$\alpha_{ij}(h_i, h_j)$. \emph{Espectralmente}, en la base de Fourier del dominio
ambas corresponden a multiplicaci\'on por una funci\'on de transferencia ---un
filtro espectral fijo $\hat{k}(\lambda)$ para la convoluci\'on, uno adaptativo
$\hat\alpha(\lambda;h)$ para la atenci\'on (cf.\
\Cref{subsec:apps:integral-operator}). \emph{Geom\'etricamente}, la estructura de
vecindarios fija el soporte de integraci\'on: la convoluci\'on integra sobre
vecindarios fijos, la atenci\'on integra con pesos que dependen del contenido
pero respetan la geometr\'ia subyacente.
\end{proof}

La \Cref{tab:apps:cnn-attn} registra la correspondencia concepto a concepto. El
mismo principio unificador ---operaciones que conmutan con la acci\'on del
grupo--- subyace a la convoluci\'on de grupo y a la atenci\'on $G$-equivariante,
a los filtros espectrales fijos y adaptativos, a la convoluci\'on \emph{gauge} y
a la atenci\'on \emph{gauge}-equivariante, y al transporte paralelo y a la
atenci\'on riemanniana. As\'i, las arquitecturas del aprendizaje profundo
geom\'etrico forman una jerarqu\'ia natural de generalizaci\'on, con la familia
convolucional r\'igida de un lado y la familia adaptativa de la atenci\'on del
otro, duales bajo esta correspondencia.

\begin{table}[h]
\centering
\caption{Correspondencias fundamentales entre la convoluci\'on geom\'etrica y la
atenci\'on geom\'etrica.}
\label{tab:apps:cnn-attn}
\begin{tabular}{@{}p{3.5cm}p{3.5cm}p{6cm}@{}}
\toprule
\textbf{CNN geom\'etrica} & \textbf{Atenci\'on geom\'etrica} & \textbf{Principio unificador} \\
\midrule
Convoluci\'on de grupo (\Cref{ch:groups}) & Atenci\'on $G$-equivariante & Operaciones que conmutan con la acci\'on de $G$ \\
\addlinespace
Filtro espectral $\hat{f}(\lambda) = g(\lambda)\hat{x}(\lambda)$ & Atenci\'on espectral $A = \softmax(f(\lambda_i,\lambda_j))$ & Multiplicaci\'on en el dominio de Fourier del laplaciano \\
\addlinespace
Convoluci\'on \emph{gauge} (\Cref{ch:gauges}) & Atenci\'on \emph{gauge}-equivariante & Respeto a las transformaciones locales de coordenadas \\
\addlinespace
\emph{Steerable} CNN, $f\mapsto\rho(g)f$ & Atenci\'on \emph{steerable}, $A\mapsto\rho(g)A\rho(g)^{-1}$ & Acci\'on coherente sobre campos vectoriales \\
\addlinespace
CNN esf\'erica, arm\'onicos $Y_\ell^m$ & Atenci\'on esf\'erica & An\'alisis arm\'onico en la esfera \\
\addlinespace
CNN sobre variedades, transporte paralelo & Atenci\'on riemanniana & Geometr\'ia intr\'inseca de las variedades \\
\bottomrule
\end{tabular}
\end{table}

\begin{remark}[Rigidez frente a adaptabilidad]
\label{rem:apps:rigidity-adaptability}
La distinci\'on m\'as fundamental entre los paradigmas convolucional y de
atenci\'on es el balance entre la estructura impuesta (rigidez) y la flexibilidad
aprendible (adaptabilidad). Cabe medir la \emph{rigidez} de una arquitectura por
la fracci\'on de su presupuesto de par\'ametros dedicada a simetr\'ias impuestas
y su \emph{adaptabilidad} por la fracci\'on dedicada a pesos dependientes del
contenido, compitiendo ambas por el mismo presupuesto. El \emph{trade-off} es
n\'itido: las garant\'ias te\'oricas escalan con la rigidez, mientras que la
capacidad expresiva escala con la adaptabilidad. A lo largo de la progresi\'on
desde la CNN cl\'asica (alta rigidez, baja adaptabilidad), pasando por las CNN de
grupo y la atenci\'on sobre grafos, hasta el \emph{transformer} completo (baja
rigidez, alta adaptabilidad), esto explica por qu\'e los modelos convolucionales
generalizan mejor con pocos datos mientras que los \emph{transformers} requieren
datos masivos pero alcanzan mayor expresividad.
\end{remark}

\begin{remark}[Convergencia de paradigmas]
\label{rem:apps:convergence}
Las arquitecturas recientes sugieren una convergencia conceptual de ambas
familias. Los dise\~nos h\'ibridos combinan un n\'ucleo convolucional ---para
caracter\'isticas locales con garant\'ias geom\'etricas--- con cabezas de
atenci\'on ---para dependencias globales adaptativas--- y un mecanismo aprendible
que las equilibra; ambas familias pueden aproximarse arbitrariamente bien por
operadores espectrales de dimensi\'on finita, y ambas convergen al mismo l\'imite
de operadores equivariantes sobre el dominio. El futuro del aprendizaje profundo
geom\'etrico no reside en la competencia entre paradigmas, sino en su s\'intesis:
combinar la estabilidad y la eficiencia de las convoluciones geom\'etricas con la
flexibilidad y la expresividad de la atenci\'on adaptativa, guiadas por
principios matem\'aticos y no por heur\'isticas. La geometr\'ia, como tantas
veces en la historia de las matem\'aticas y la f\'isica, proporciona el lenguaje
unificador.
\end{remark}

\section{Visi\'on por computador}
\label{sec:apps:vision}

La visi\'on por computador es la cuna hist\'orica de las ideas de este libro.
Las redes neuronales convolucionales ganaron la competici\'on ImageNet en
2012~\citep{krizhevsky2012alexnet} precisamente porque codifican una simetr\'ia
de las im\'agenes: una traslaci\'on de la entrada produce un mapa de
caracter\'isticas trasladado. Como se desarrolla en \Cref{ch:grids}, esta
equivarianza a traslaciones es lo que permite reutilizar un \'unico filtro
aprendido en cada posici\'on, reduciendo dr\'asticamente el n\'umero de
par\'ametros y proporcionando el sesgo inductivo que hace que las im\'agenes
sean aprendibles a partir de datos finitos. Todo lo que sigue en visi\'on por
computador puede leerse como un intento de extender este \'exito m\'as all\'a de
la rejilla plana de p\'ixeles.

\paragraph{M\'as all\'a de la traslaci\'on.}
Muchas tareas visuales son invariantes a m\'as que la traslaci\'on. Una galaxia,
una c\'elula al microscopio o una imagen a\'erea no tienen una orientaci\'on
preferida; un clasificador no deber\'ia tener que reaprender el mismo patr\'on en
cada copia rotada. Las redes convolucionales equivariantes a grupos
(\Cref{ch:groups}) incorporan la simetr\'ia de rotaci\'on y reflexi\'on
directamente en la convoluci\'on, de modo que una caracter\'istica detectada en
una orientaci\'on se detecta autom\'aticamente en todas. El beneficio es mayor en
conjuntos peque\~nos, donde el aumento de datos no cubre el grupo de simetr\'ias
con suficiente densidad y la arquitectura ha de aportar la invarianza.

\paragraph{Visi\'on tridimensional.}
La ruptura m\'as limpia con la rejilla llega con los datos tridimensionales. Un
escaneo LiDAR o un sensor de profundidad devuelve una \emph{nube de puntos}: un
conjunto desordenado de puntos en $\R^3$, sin rejilla, sin orden can\'onico y de
cardinalidad variable. La simetr\'ia relevante es doble. Primero, la salida debe
ser \emph{invariante a las permutaciones} de los puntos, pues el conjunto
$\{p_1,\dots,p_n\}$ no depende de c\'omo enumeremos sus elementos. Segundo, en
muchas tareas debe ser invariante o equivariante a los movimientos r\'igidos de
la escena. El primer requisito lo capta exactamente la teor\'ia de funciones
sobre conjuntos.

\begin{theorem}[Representaci\'on universal de funciones de conjuntos]
\label{thm:apps:deepsets}
Toda funci\'on continua e invariante a permutaciones $f$ de un conjunto
$\{x_1,\dots,x_n\}\subset\R^d$ puede escribirse en la forma
\[
  f(\{x_1,\dots,x_n\}) \;=\; \rho\!\left(\sum_{i=1}^n \phi(x_i)\right),
\]
para funciones continuas $\phi$ y $\rho$ apropiadas.
\end{theorem}

Este resultado, debido al marco DeepSets~\citep{zaheer2017deepsets}, es el
fundamento te\'orico de arquitecturas como PointNet~\citep{qi2017pointnet}:
codificar cada punto de forma independiente con una red compartida $\phi$,
agrupar las codificaciones con un agregador invariante a permutaciones (una suma
o un m\'aximo) y decodificar con $\rho$. La construcci\'on es la instancia m\'as
sencilla posible del esquema de paso de mensajes de \Cref{ch:graphs} ---una
\'unica ronda de mensajes en un grafo sin aristas--- y ya garantiza la
invarianza correcta por dise\~no en lugar de por aumento de datos.

\paragraph{Dominios curvos.}
Cuando los datos viven sobre una superficie curva en lugar de en el espacio
libre, toman el relevo los temas del \emph{gauge} y de las geod\'esicas
(\Cref{ch:geodesics,ch:gauges}). Las se\~nales esf\'ericas ---im\'agenes
panor\'amicas, mapas cosmol\'ogicos, la superficie de la Tierra--- se procesan
con redes convolucionales esf\'ericas~\citep{cohen2018sphericalcnns}
equivariantes al grupo de rotaciones $\SO(3)$, de modo que rotar la esfera rota
coherentemente la respuesta. Sobre una malla general no hay un marco global, y
las redes \emph{gauge}-equivariantes garantizan que la salida no dependa de la
elecci\'on local arbitraria de coordenadas. En todos los casos la lecci\'on
recurrente es la misma: identificar la simetr\'ia del dominio y construirla en el
operador.

\section{Mol\'eculas y sistemas f\'isicos}
\label{sec:apps:molecules}

Si hay un dominio que vindica el aprendizaje profundo geom\'etrico, es el de las
ciencias f\'isicas y de la vida. Las mol\'eculas \emph{son} grafos cuyos nodos
son \'atomos y cuyas aristas son enlaces, embebidos en $\R^3$; las leyes que las
gobiernan son manifiestamente invariantes bajo los movimientos r\'igidos del
espacio. Un modelo que ignore estas simetr\'ias malgasta capacidad reaprendiendo
una f\'isica que podr\'ia simplemente asumirse.

\paragraph{Predicci\'on de propiedades moleculares.}
Las redes neuronales de grafos predicen propiedades farmacol\'ogicas y
cu\'antico-qu\'imicas directamente a partir del grafo molecular, reemplazando los
descriptores artesanales y las codificaciones en cadena (como SMILES) de la
quimioinform\'atica cl\'asica, que descartan la informaci\'on tridimensional que
gobierna la actividad biol\'ogica. Cuando el objetivo depende de la
conformaci\'on molecular ---energ\'ias, fuerzas, momentos dipolares--- el modelo
debe ser $\Eucl(3)$- o $\SE(3)$-equivariante: rotar la mol\'ecula debe rotar los
vectores predichos y dejar inalterados los escalares predichos. Las arquitecturas
de atenci\'on equivariante de la \Cref{subsec:apps:attn-mpnn}, junto con las redes
de campos tensoriales y de paso de mensajes
equivariante~\citep{satorras2021enequivariant}, son precisamente las
herramientas que imponen esto.

\begin{example}[AlphaFold y la estructura de las prote\'inas]
\label{ex:apps:alphafold}
El \'exito m\'as celebrado del campo es
AlphaFold~\citep{jumper2021alphafold}, que resolvi\'o el problema, abierto
durante d\'ecadas, de predecir la estructura tridimensional de una prote\'ina a
partir de su secuencia de amino\'acidos. En su n\'ucleo hay una arquitectura
basada en atenci\'on que razona sobre la geometr\'ia de los residuos en el
espacio respetando la simetr\'ia $\SE(3)$ de los movimientos r\'igidos: la
estructura predicha se transforma de forma consistente cuando el marco de
referencia se rota o se traslada. La equivarianza geom\'etrica no es aqu\'i una
elecci\'on estil\'istica sino un requisito estructural, y es en buena medida la
raz\'on por la que el modelo generaliza tan bien. AlphaFold es, en este sentido,
la ilustraci\'on paradigm\'atica de la conexi\'on teor\'ia--arquitectura que este
libro traza: la equivarianza $\SE(3)$ abstracta de \Cref{ch:groups} hecha
concreta en un sistema que transform\'o la biolog\'ia estructural.
\end{example}

\paragraph{Simulaci\'on f\'isica.}
Los mismos principios se extienden de las estructuras est\'aticas a la
din\'amica. Las redes equivariantes aprenden a simular sistemas f\'isicos
---fluidos, materiales deformables, din\'amica de muchos cuerpos--- operando
sobre mallas irregulares o conjuntos de part\'iculas mientras respetan las
simetr\'ias de las ecuaciones diferenciales subyacentes. Como la simetr\'ia est\'a
incorporada, las predicciones son f\'isicamente consistentes por construcci\'on
(una entrada rotada produce una predicci\'on correspondientemente rotada) y los
modelos suelen igualar o superar a los solucionadores num\'ericos cl\'asicos a una
fracci\'on del coste. Un \'unico modelo fundacional puede incluso preentrenarse a
la vez sobre mol\'eculas peque\~nas y prote\'inas: la equivarianza es una
propiedad arquitect\'onica, independiente de los datos, de modo que una red
entrenada para ser equivariante en un dominio sigue si\'endolo al transferirse a
otro con la misma simetr\'ia~\citep{jiao2024ept}.

\paragraph{Ciencia de materiales.}
Los s\'olidos cristalinos a\~naden una capa adicional de simetr\'ia: las
simetr\'ias discretas de traslaci\'on y de grupo puntual de la red cristalina.
Codificar estas simetr\'ias en redes equivariantes permite a los modelos predecir
propiedades mec\'anicas, t\'ermicas y electr\'onicas de compuestos candidatos,
acelerando la b\'usqueda de nuevos materiales. Como en el caso molecular, el
\emph{prior} geom\'etrico convierte etiquetas escasas y costosas en un problema
de aprendizaje tratable.

\section{Sistemas de recomendaci\'on y redes sociales}
\label{sec:apps:graphs}

Los datos relacionales ---qui\'en interact\'ua con qui\'en, qui\'en compra
qu\'e--- tienen forma de grafo por naturaleza, y aqu\'i la simetr\'ia que rige es
la invarianza a permutaciones: la identidad de un usuario no depende del orden en
que enumeremos las cuentas. Las redes neuronales de grafos (\Cref{ch:graphs})
explotan esto directamente, propagando informaci\'on a lo largo de las aristas
del grafo de interacciones de modo que la representaci\'on de cada nodo queda
modelada por su vecindario.

\paragraph{Recomendaci\'on y an\'alisis social.}
Los sistemas de recomendaci\'on modelan usuarios e \'items como nodos de un grafo
bipartito y aprenden \emph{embeddings} mediante paso de mensajes sobre las
interacciones observadas, capturando se\~nales colaborativas que las tablas
planas de caracter\'isticas pasan por alto. El an\'alisis de redes sociales usa la
misma maquinaria para la clasificaci\'on de nodos (inferir atributos de los
usuarios), la predicci\'on de enlaces (sugerir conexiones) y la detecci\'on de
comunidades. En todos estos casos la topolog\'ia irregular ---nodos con n\'umeros
de vecinos muy dispares y sin orden natural--- es exactamente lo que los modelos
convolucionales no pueden manejar y para lo que las redes de grafos est\'an
construidas.

\paragraph{Jerarqu\'ias y geometr\'ia hiperb\'olica.}
Muchos conjuntos de datos relacionales no son solo irregulares sino
\emph{jer\'arquicos}: taxonom\'ias, organigramas, grafos de conocimiento y la
estructura de seguidores de las redes sociales se ramifican como \'arboles. El
n\'umero de nodos a distancia $r$ de una ra\'iz crece exponencialmente en $r$, y
el espacio euclidiano simplemente no tiene sitio suficiente para embeber tal
crecimiento sin aplastar las distancias. El espacio hiperb\'olico s\'i lo tiene:
su volumen tambi\'en crece exponencialmente, de modo que un \'arbol se embebe con
casi ninguna distorsi\'on (\Cref{fig:apps:hyperbolic}). Esta es la motivaci\'on
geom\'etrica detr\'as de los \emph{embeddings} hiperb\'olicos y de las redes de
atenci\'on hiperb\'olica de la \Cref{subsec:apps:attn-mpnn}, que sit\'uan las
entidades de los grafos de conocimiento y los usuarios jer\'arquicos en un
espacio de curvatura negativa.

\begin{figure}[h]
\centering
\begin{tikzpicture}[
    root node/.style={circle, draw=gdlpurple!80, fill=gdlpurple!70,
        minimum size=12pt, inner sep=0pt},
    depth1 node/.style={circle, draw=gdlblue!80, fill=gdlblue!60,
        minimum size=10pt, inner sep=0pt},
    depth2 node/.style={circle, draw=gdlgreen!80, fill=gdlgreen!60,
        minimum size=10pt, inner sep=0pt},
    leaf node/.style={circle, draw=gdlyellow!80, fill=gdlyellow!70,
        minimum size=8pt, inner sep=0pt},
    tree edge/.style={draw=gdlgray!50, line width=0.8pt},
    scale=0.82, every node/.style={transform shape}
]
% IZQUIERDA: euclidiano
\begin{scope}[shift={(0,0)}]
    \coordinate (r) at (0, 1.9);
    \coordinate (d1-0) at (-1.8, 0.4); \coordinate (d1-1) at ( 1.8, 0.4);
    \coordinate (d2-0) at (-2.7, -1.1); \coordinate (d2-1) at (-0.9, -1.1);
    \coordinate (d2-2) at ( 0.9, -1.1); \coordinate (d2-3) at ( 2.7, -1.1);
    \coordinate (d3-0) at (-3.15, -2.6); \coordinate (d3-1) at (-2.25, -2.6);
    \coordinate (d3-2) at (-1.35, -2.6); \coordinate (d3-3) at (-0.45, -2.6);
    \coordinate (d3-4) at ( 0.45, -2.6); \coordinate (d3-5) at ( 1.35, -2.6);
    \coordinate (d3-6) at ( 2.25, -2.6); \coordinate (d3-7) at ( 3.15, -2.6);
    \draw[dashed, gdlgray!50, line width=0.8pt] (0,-0.35) circle (3.6);
    \draw[tree edge] (r) -- (d1-0); \draw[tree edge] (r) -- (d1-1);
    \draw[tree edge] (d1-0) -- (d2-0); \draw[tree edge] (d1-0) -- (d2-1);
    \draw[tree edge] (d1-1) -- (d2-2); \draw[tree edge] (d1-1) -- (d2-3);
    \draw[tree edge] (d2-0) -- (d3-0); \draw[tree edge] (d2-0) -- (d3-1);
    \draw[tree edge] (d2-1) -- (d3-2); \draw[tree edge] (d2-1) -- (d3-3);
    \draw[tree edge] (d2-2) -- (d3-4); \draw[tree edge] (d2-2) -- (d3-5);
    \draw[tree edge] (d2-3) -- (d3-6); \draw[tree edge] (d2-3) -- (d3-7);
    \node[root node] at (r) {};
    \node[depth1 node] at (d1-0) {}; \node[depth1 node] at (d1-1) {};
    \foreach \i in {0,...,3} {\node[depth2 node] at (d2-\i) {};}
    \foreach \i in {0,...,7} {\node[leaf node] at (d3-\i) {};}
    \node[font=\footnotesize\itshape, text=gdlgray!70!black] at (0, -4.6)
        {Embebido euclidiano};
\end{scope}
% DERECHA: disco de Poincar\'e
\begin{scope}[shift={(9.0, -0.35)}]
    \pgfmathsetmacro{\R}{3.6}
    \foreach \rad in {0.30, 0.52, 0.70, 0.83, 0.92} {
        \draw[gdlgray!20, thin] (0,0) circle (\rad*\R);}
    \foreach \ang in {0,30,60,90,120,150,180,210,240,270,300,330} {
        \draw[gdlgray!20, thin] (0,0) -- ({\ang}:\R);}
    \draw[black, line width=1.2pt] (0,0) circle (\R);
    \pgfmathsetmacro{\rone}{0.6044*\R}
    \pgfmathsetmacro{\rtwo}{0.8854*\R}
    \pgfmathsetmacro{\rthree}{0.9505*\R}
    \coordinate (pr) at (0, 0);
    \coordinate (pd1-0) at ({90}:{\rone}); \coordinate (pd1-1) at ({270}:{\rone});
    \coordinate (pd2-0) at ({50}:{\rtwo}); \coordinate (pd2-1) at ({130}:{\rtwo});
    \coordinate (pd2-2) at ({230}:{\rtwo}); \coordinate (pd2-3) at ({310}:{\rtwo});
    \coordinate (pd3-0) at ({35}:{\rthree}); \coordinate (pd3-1) at ({65}:{\rthree});
    \coordinate (pd3-2) at ({115}:{\rthree}); \coordinate (pd3-3) at ({145}:{\rthree});
    \coordinate (pd3-4) at ({215}:{\rthree}); \coordinate (pd3-5) at ({245}:{\rthree});
    \coordinate (pd3-6) at ({295}:{\rthree}); \coordinate (pd3-7) at ({325}:{\rthree});
    \draw[tree edge] (pr) -- (pd1-0); \draw[tree edge] (pr) -- (pd1-1);
    \draw[tree edge] (pd1-0) -- (pd2-0); \draw[tree edge] (pd1-0) -- (pd2-1);
    \draw[tree edge] (pd1-1) -- (pd2-2); \draw[tree edge] (pd1-1) -- (pd2-3);
    \draw[tree edge] (pd2-0) -- (pd3-0); \draw[tree edge] (pd2-0) -- (pd3-1);
    \draw[tree edge] (pd2-1) -- (pd3-2); \draw[tree edge] (pd2-1) -- (pd3-3);
    \draw[tree edge] (pd2-2) -- (pd3-4); \draw[tree edge] (pd2-2) -- (pd3-5);
    \draw[tree edge] (pd2-3) -- (pd3-6); \draw[tree edge] (pd2-3) -- (pd3-7);
    \node[root node] at (pr) {};
    \node[depth1 node] at (pd1-0) {}; \node[depth1 node] at (pd1-1) {};
    \foreach \i in {0,...,3} {\node[depth2 node] at (pd2-\i) {};}
    \foreach \i in {0,...,7} {\node[leaf node] at (pd3-\i) {};}
    \node[font=\footnotesize\itshape, text=gdlgray!70!black] at (0, -4.25)
        {Embebido en el disco de Poincar\'e};
\end{scope}
\end{tikzpicture}
\caption[Embebido euclidiano frente a hiperb\'olico de un \'arbol]{Embebido de un
\'arbol equilibrado en el espacio euclidiano (izquierda) y en el disco de
Poincar\'e (derecha). En el plano, los niveles sucesivos se amontonan y los
sub\'arboles hermanos se solapan, distorsionando las distancias; en el espacio
hiperb\'olico el volumen que crece exponencialmente da a cada nivel sitio para
extenderse, de modo que el \'arbol se embebe con poca distorsi\'on. Por eso los
datos relacionales jer\'arquicos se modelan de forma natural en un espacio de
curvatura negativa.}
\label{fig:apps:hyperbolic}
\end{figure}

\paragraph{Grafos temporales: fraude y riesgo.}
Las transacciones financieras forman grafos cuya topolog\'ia evoluciona en el
tiempo. Las redes neuronales de grafos que operan sobre estos grafos temporales
detectan redes de fraude y riesgo sist\'emico reconociendo patrones estructurales
---subgrafos inusuales, cambios bruscos en la conectividad--- que los modelos
tabulares, ciegos a la estructura relacional, no pueden ver. El \emph{prior}
geom\'etrico es de nuevo decisivo: es la \emph{forma} del patr\'on de
interacci\'on, y no las caracter\'isticas individuales, la que porta la se\~nal.

\section{Un principio com\'un}
\label{sec:apps:common}

A lo largo de estos dominios reaparece un \'unico mensaje. En visi\'on por
computador la simetr\'ia es la traslaci\'on, la rotaci\'on o la libertad
\emph{gauge} de una superficie curva; en el modelado molecular y f\'isico es el
grupo euclidiano de los movimientos r\'igidos; en los datos relacionales es la
permutaci\'on de nodos, y la geometr\'ia del espacio de embebido puede ser plana
o de curvatura negativa. En todos los casos se aplica la misma receta:
identificar la simetr\'ia que el problema respeta, construirla en la arquitectura
y dejar que la red gaste su capacidad en lo que realmente var\'ia.

\begin{remark}[La geometr\'ia como principio organizador]
\label{rem:apps:unification}
La amplitud de estas aplicaciones es, en s\'i misma, la mayor evidencia de la
tesis central de este libro. Las redes convolucionales para im\'agenes, los
\emph{transformers} equivariantes para prote\'inas y las redes de grafos para
datos sociales parecen herramientas inconexas, y sin embargo cada una es una
instancia del paso de mensajes equivariante, que difieren \'unicamente en el
dominio y el grupo de simetr\'ias que respetan ---y, por la
\Cref{thm:apps:cnn-attn-correspondence}, que difieren solo en si los mensajes
portan pesos fijos o adaptativos. La geometr\'ia no es un truco espec\'ifico de
cada dominio, sino el principio com\'un que el aprendizaje profundo geom\'etrico
identifica, formaliza y explota.
\end{remark}

M\'as all\'a de estas \'areas consolidadas, las mismas ideas abren l\'ineas de
investigaci\'on dondequiera que la estructura geom\'etrica sea fundamental: la
neurociencia computacional, donde las redes de grafos modelan la conectividad
cerebral estructural y funcional y revelan patrones de trastorno neurol\'ogico
invisibles para los m\'etodos tabulares; y la climatolog\'ia, donde la geometr\'ia
esf\'erica de la Tierra invita a redes esf\'ericas $\SO(3)$-equivariantes para la
predicci\'on meteorol\'ogica y clim\'atica. La diversidad de estos escenarios
subraya que respetar la geometr\'ia de los datos no es un detalle de
optimizaci\'on, sino el principio organizador del campo, y el puente desde la
teor\'ia de los cap\'itulos precedentes hasta los sistemas que est\'an
remodelando la ciencia y la tecnolog\'ia.
